\chapter{整数下降的准备}\label{sec:descent}

\section{为什么选择整数下降}
设 $E(x)=\sum_i x_i^2$。一次平均的精确能量差是
\begin{align*}
E(x)-E(T_Sx)
 &=\sum_{i\in S}x_i^2-p\bar x_S^2\\
 &=\sum_{i\in S}(x_i-\bar x_S)^2.
\end{align*}
因此，只要所选数不全相等，能量就严格下降。
但在有理数上，严格下降的正数列仍可能无限长。
本部分的策略是让每次平均都保持整数；这样 $E$ 始终为非负整数，下降必然有限。

要使这条策略成立，还必须避免平均过程制造新的模素数障碍。
我们将同时维持四项条件：整数、零和、每个危险素数处不全同余、
至少两个不同数值各占据 $p$ 个位置。
最后一项是后续选择的储备：每次既能产生一个新的 $p$ 重值，又能留住一个旧的 $p$ 重值。
若已处于明确的相反值终端，就直接消去，不再要求维持这项储备。

\section{危险集与未选见证}

本讲设 $p\ge5$，$d=d_p(n)$。称重数至少为 $p$ 的实际数值为重值。构造维持整数、零和、合法以及至少两个不同重值。


\begin{definition}[危险素数与危险集]
称 $\ell\mid n$ 且 $\ell\ne p$ 为危险素数。
若模 $\ell$ 有一个余数类包含至少 $n-p$ 个位置，称它为主类，
将其余位置构成的集合记为 $F_\ell$。
如果一次平均会使全部位置同余，那么未选的 $n-p$ 个位置已经必须同余。
因此只有这种主类存在时，才有一步制造障碍的风险。
\end{definition}
有些危险素数并无这样的主类；它们对此次整数选择没有额外限制。
有主类时，只需让一个主类外的位置留在选择之外。
而由于主类大小大于 $p$，任意选择后也会留下主类位置。
这就是“预留一个见证”的确切含义。

\begin{lemma}[由未选位置保证安全]\label{lem:safety}
设 $n>2p$。若模危险素数 $\ell$ 存在大小至少 $n-p$ 的余类，其补集 $F_\ell$ 唯一，且 $2\le|F_\ell|\le p$。不包含任何完整 $F_\ell$ 的整数平均安全。以下整数选择也安全：
\begin{enumerate}
\item $(a^{p-1},x)$，其中至少一份 $a$ 未选；
\item 由若干份 $a,b$ 及一份第三值 $x$ 组成，且 $a,b$ 均有未选副本；
\item 每种被选数值均有未选副本。
\end{enumerate}
\end{lemma}
\begin{proof}
$n-p>n/2$ 给出主类唯一性。例外数不能为 $1$，否则零和模 $\ell$ 迫使唯一例外仍属主类。避开完整 $F_\ell$ 时，未选位置保留主类及例外。对后三种选择，若输出全同余，未选副本及选中和为 $p$ 倍新均值依次迫使全部被选值也属于共同余类，矛盾。
\end{proof}


\section{为什么两步就能建立储备}
初始时可以把分析代表由 $X$ 换成 $pX$。
这样所有初值都被 $p$ 整除，任意 $p$ 项的第一次平均都是整数。
我们不把乘 $p$ 作为新增操作；归一化引理说明，对 $pX$ 找到的同一平均词也适用于 $X$。
后面的所有下降都固定在这个整数尺度内进行。

初始化中经常使用一个简单的“删一”观察：
对 $p+1$ 个数 $x_1,\ldots,x_{p+1}$，删去第 $i$ 个数后的和为 $S-x_i$。
若这些 $p+1$ 个候选和都相同，则全部 $x_i$ 相同。
因而，只要池中没有某实际值的 $p$ 份副本，就能用多出的一个位置避开一个指定均值。

\begin{lemma}[建立两个不同重值]\label{lem:initialize}
若 $n\ge3p+d$，从 $pX$ 出发至多两次整数平均可建立两个不同重值。
\end{lemma}
\begin{proof}
在任一候选池中，为仍完整包含在池内的每个 $F_\ell$ 留出一个位置，至多留出 $d$ 个。若池内所有值同余模 $p$、每种实际值不足 $p$ 份，余下 $p$ 个位置给安全非恒等整数平均；余下 $p+1$ 个位置时还可避开一个指定均值，因为所有删一子集同和会迫使这 $p+1$ 个值全相等。
已有两个重值时无须操作。仅一个重值 $a$ 且重数至少 $2p-1$ 时，选择 $(a^{p-1},b)$，留下 $p$ 份 $a$ 并生成另一重值。若重数至多 $2p-2$，在 $a$ 外的池中避开均值 $a$。若初始没有重值，先制造一个；它的重数至多 $2p-1$，其外至少 $d+p+1$ 个未改变的初值仍被 $p$ 整除，故第二次可制造不同重值。
\end{proof}


\section{模p零和选择：足够的池与唯一的边界形状}
整性要求所选 $p$ 项之和为 $0\pmod p$。
接下来的零和引理告诉我们，$2p-1$ 个候选位置一定足够；
若只能拿出 $2p-2$ 个位置，失败的余数形状也完全确定。
这两种结论将分别用于宽裕的选择和最后一个容量边界。

\begin{example}[五平均的边界池]
模 $5$ 的池 $(0^4,1^4)$ 没有五项零和。
五项中必须包含 $1,2,3$ 或 $4$ 个 $1$，和都不被 $5$ 整除。
添加任意一个第九位置以后，零和引理保证出现五项零和。
逆命题说，长度为八的失败池在仿射改名后只能有这种两类各四份的形状。
\end{example}

\begin{lemma}[素数零和选择及逆边界]\label{lem:egz}
任意 $2p-1$ 个 $\F_p$ 元素包含 $p$ 项零和。长度 $2p-2$ 的序列没有 $p$ 项零和，当且仅当它为 $u^{p-1}v^{p-1}$、$u\ne v$。
\end{lemma}
\begin{proof}
先说明有限域计数：若 $m$ 元多项式 $F_i$ 的次数和小于 $m$，其共同零点数模 $p$ 等于
\[\sum_{z\in\F_p^m}\prod_i(1-F_i(z)^{p-1})=0.
\]
这里把 $1-F_i(z)^{p-1}$ 看成一个指示函数：$F_i(z)=0$ 时它为 $1$，否则为 $0$。
对一个变量，有 $\sum_{t\in\F_p}t^k=0$，除非 $k$ 是正的 $p-1$ 倍数；当 $k=0$ 时求和也是 $p=0$。
展开后每个单项式的总次数小于 $m(p-1)$，至少一个变量的指数不是正的 $p-1$ 倍数，对它求和为零。应用于
\[\sum_{i=1}^{2p-1}z_i^{p-1}=0,\qquad
\sum_{i=1}^{2p-1}a_i z_i^{p-1}=0.
\]
由于 $z_i^{p-1}$ 在 $z_i\ne0$ 时为 $1$，第一个方程要求非零坐标个数被 $p$ 整除。
一共只有 $2p-1$ 个变量，所以非零解的支持恰有 $p$ 项；第二个方程要求该支持对应的 $a_i$ 之和为零。
每个合格支持允许每一非零坐标任取 $p-1$ 个值，零点数为 $1+(p-1)^p M$，故 $M\equiv1\pmod p$。

证明逆命题。出现 $p$ 次同一余数时已有零和；仅两类时只能为所示形式。若至少三类且各重数至多 $p-1$，可选两对不交位置，使两对的差非零且不互为正负。取重数最大的类 $a$，它至少有两份；若两个其他类对 $a$ 的差不互为正负即可。否则仅有 $a,b,c$ 三类，且 $b,c$ 关于 $a$ 对称；$b,c$ 至少一类有两份，改用 $(a,b),(b,c)$，两差相差倍率 $2$。
余下 $2p-6$ 项最大重数至多 $p-2$，可再配 $p-4$ 对不同余数，留下两个单点。所用匹配事实为：多重数序列的异类最大匹配数是 $\min(\lfloor M/2\rfloor,M-f_{\max})$。若某类超过一半，将其与其余类配对；否则将各类连续排成一列，前后半列配对即达到该界。
前两对各取一项的和集有四个元素。以后每加一对，和集至少增大一，除非已为 $\F_p$：一个真子集不能在非零平移下不变。故 $p-2$ 对的选择和遍历 $\F_p$，加上两个单点得到 $p$ 项零和。
\end{proof}

